\documentclass[../main.tex]{subfiles}
\begin{document}
\begin{center}
    \Large{\textbf{ABSTRACT}}\\
\end{center}
\vspace{1cm}
\indent Food waste and hunger form a coexisting paradox in Vietnam: a large amount of still-edible food is discarded every day, while many disadvantaged individuals and organizations remain short of meals. Food-sharing activities today are mostly spontaneous and fragmented, carried out through social networks or phone calls, which makes it hard for donors, recipients, and intermediate carriers to connect in time, especially given how perishable food is.

\indent To address this problem, the project develops Food 4 life, a real-time platform for sharing surplus food tailored to the Vietnamese context. The system follows a client--server architecture and is built with a cross-platform mobile application (React Native/Expo), a Node.js/Express server, a MongoDB database, and a Socket.IO real-time channel.

\indent The system supports four roles --- donor, recipient, volunteer, and administrator --- and enables a two-way workflow in which donors post surplus food while recipients can proactively post requests. It also integrates location-based suggestions, semi-automatic volunteer dispatch, in-app messaging, multilingual push notifications, and delivery-status tracking.

\indent The main contributions are a set of technical solutions that ensure correctness under non-ideal operating conditions: safe matching under concurrency, a hybrid distance estimator with graceful fallback, a self-healing order state machine, an anti-fraud pickup code, and instantly revocable stateless authentication. As a result, the project delivers a complete product consisting of a mobile application for three user roles, a web-based administration portal, and a backend server, all operating stably across the main business flows.
\end{document}